% File: manual.tex
% Author: Y.U.P. (yashparitkar)
% Last Modified: 2026-08-10 Mon 20:29
%
% Description: LaTeX source file for the LibreILA user manual. This covers mainly the usage instructions for the LibreILA tool.

\documentclass{article}
\usepackage{float}
\usepackage[margin=2.5cm]{geometry}
% \usepackage{xurl}
\usepackage{tabularx}
\usepackage[utf8x]{inputenc}
\usepackage{amssymb}
\usepackage{amsmath}
\usepackage{xcolor}
\usepackage{colortbl}
\usepackage[linktocpage]{hyperref}
% \usepackage{wasysym}
\usepackage{graphicx}
% \usepackage{epsfig}
\setkeys{Gin}{width=0.8\textwidth}
% \usepackage[normalem]{ulem}
% \usepackage{enumerate}
\usepackage{listings}
\usepackage[nohyperlinks]{acronym}

\lstset{ %
basicstyle=\footnotesize,       % the size of the fonts that are used for the code
showspaces=false,               % show spaces adding particular underscores
showstringspaces=false,         % underline spaces within strings
showtabs=false,                 % show tabs within strings adding particular underscores
frame=single,                   % adds a frame around the code
tabsize=2,                      % sets default tabsize to 2 spaces
breaklines=true,                % sets automatic line breaking
breakatwhitespace=false,        % sets if automatic breaks should only happen at whitespace
}

\hypersetup{
 colorlinks=true,
 citecolor=violet,
 linkcolor=blue,
 urlcolor=blue
}

\usepackage{fancyhdr}


\fancyhead[L]{LibreILA User Manual}
\fancyhead[R]{Yash Paritkar}
\fancyfoot[L]{\leftmark}
\fancyfoot[C]{}
\fancyfoot[R]{\thepage}
\renewcommand{\headrulewidth}{0.5pt}
\renewcommand{\footrulewidth}{0.5pt}
\pagestyle{fancy}


\newcommand{\todo}[1]{\textbf{\color{red}TODO: #1}}
\newcommand{\arya}{\color{blue}} 
\newcommand{\degr}{\ensuremath{^\circ}}

\title{LibreILA User Manual}
\author{Yash Paritkar (github/yashparitkar)}
\date{July 2026}

\begin{document}

\begin{figure}
    \centering
    \includegraphics[trim=0 0 0 0,clip=true,width=0.75\linewidth]{./images/libre_ila_logo_converted.pdf}
    % \caption{Caption}
    \label{fig:rendering}
\end{figure}

\maketitle

\newpage

\section*{Document History}
\addcontentsline{toc}{section}{Document history}

\renewcommand{\arraystretch}{1.5}
\begin{table}[H]
    \centering
    \begin{tabularx}{\textwidth}{|l|c|X|}
        \hline
        \rowcolor{gray!40} \textbf{Version} & \textbf{Date} & \textbf{Comments} \\
        \hline
        0.0 & 25 July 2026 & Document created.\\
        \hline
        \end{tabularx}
%    \caption{General characteristics of the payload}
    \label{tab:hisotry}
\end{table}

\newpage

\tableofcontents

\newpage
\listoffigures
\listoftables

\section*{List of Acronyms}
\begin{acronym}[sort = true]
    % A
    \acro{AXI}{Advanced eXtensible Interface}
    % C
    \acro{CDC}{Clock Domain Crossing}
    % F
    \acro{FIFO}{First In First Out (queue)}
    \acro{FPGA}{Field Programmable Gate Array}
    % G
    \acro{GUI}{Graphical User Interface}
    % H
    \acro{HDL}{Hardware Description Language}
    % I
    \acro{ILA}{Integrated Logic Analyzer}
    \acro{IP}{Intellectual Property}
    % L
    \acro{LSRAM}{Large Static Random Access Memory}
    % O
    \acro{OHL}{Open Hardware Licence}
    % R
    \acro{RTL}{Register Transfer Level}
    % S
    \acro{SDC}{Synopsys Design Constraints}
    % U
    \acro{UART}{Universal Asynchronous Receiver Transmitter}
    % V
    \acro{VCD}{Value Change Dump}
\end{acronym}

\newpage
\section{Introduction}
This document describes the design and implementation of the \ac{ILA}.

\newpage
\section{Generating configuration}
\begin{abstract}
    If same sampling clock speed is ok for multiple channels, a single ILA can be used to sample multiple channels otherwise multiple ILAs can be used to sample different channels at different speeds.
\end{abstract}

\newpage
\section{Creating a generic logic analyzer from LibreILA}

\newpage
\section{Using the CLI interface of LibreILA}

\newpage
\section{Using the GUI interface of LibreILA}
\end{document}
